\documentclass[11pt,a4paper]{article}
\usepackage[utf8]{inputenc}
\usepackage[T1]{fontenc}
\usepackage{geometry}
\geometry{margin=1in}
\usepackage{hyperref}
\usepackage{enumitem}
\usepackage{booktabs}
\usepackage{longtable}
\usepackage{xcolor}
\usepackage{titlesec}
\usepackage{amsmath,amssymb}

\titleformat{\section}{\Large\bfseries\color{blue!70!black}}{}{0em}{}
\titleformat{\subsection}{\large\bfseries\color{blue!50!black}}{}{0em}{}
\titleformat{\subsubsection}{\normalsize\bfseries\color{blue!30!black}}{}{0em}{}

\title{\textbf{Replication Plan}\\[0.5em]
\large Variational Monte Carlo Calculations of $A \leq 4$ Nuclei with an Artificial Neural-Network Correlator Ansatz\\[0.3em]
\normalsize OSTI ID: 1864334}
\author{Auto-generated Replication Blueprint}
\date{\today}

\begin{document}
\maketitle
\tableofcontents
\newpage

%---------------------------------------------------------------------
\section{Paper Overview}
%---------------------------------------------------------------------
This paper applies variational Monte Carlo (VMC) with a neural-network-based wave function ansatz to compute ground-state energies of light nuclei ($A \leq 4$: deuteron, triton, helium-3, helium-4). The neural network serves as a correlator that modifies a Slater-determinant or Jastrow baseline, parameterizing many-body correlations. The network is optimized using the stochastic reconfiguration (SR) algorithm to minimize the variational energy. The nuclear Hamiltonian uses realistic nucleon--nucleon (NN) potentials (e.g., AV6$'$, AV8$'$, or Minnesota potential).

%---------------------------------------------------------------------
\section{Detailed Workflow}
%---------------------------------------------------------------------

\subsection{Phase 1: Nuclear Hamiltonian}
\begin{enumerate}[label=\textbf{1.\arabic*}]
  \item \textbf{Implement the NN potential.}
    \begin{itemize}
      \item Central, spin--spin, tensor, spin--orbit components.
      \item Potential parameters from published tables (Argonne AV6$'$/AV8$'$ or Minnesota).
    \end{itemize}
  \item \textbf{Kinetic energy operator.}
    \begin{itemize}
      \item Non-relativistic kinetic energy $T = -\frac{\hbar^2}{2m_N} \sum_i \nabla_i^2$.
      \item Remove center-of-mass kinetic energy.
    \end{itemize}
\end{enumerate}

\subsection{Phase 2: Wave Function Ansatz}
\begin{enumerate}[label=\textbf{2.\arabic*}]
  \item \textbf{Construct the baseline wave function.}
    \begin{itemize}
      \item Slater determinant of single-particle orbitals (Gaussians or harmonic oscillator).
      \item Jastrow correlation factor for short-range pair correlations.
    \end{itemize}
  \item \textbf{Neural network correlator.}
    \begin{itemize}
      \item Input: inter-particle distances $r_{ij}$ and spin/isospin quantum numbers.
      \item Architecture: fully connected network with 2--4 hidden layers; tanh or ReLU activations.
      \item Output: multiplicative correction to the wave function $\Psi = \text{NN}(\{r_{ij}\}) \times \Psi_{\text{baseline}}$.
      \item Ensure antisymmetry under particle exchange (via the Slater determinant).
    \end{itemize}
\end{enumerate}

\subsection{Phase 3: VMC Sampling}
\begin{enumerate}[label=\textbf{3.\arabic*}]
  \item \textbf{Metropolis--Hastings sampling.}
    \begin{itemize}
      \item Sample nucleon configurations $\mathbf{R} = \{r_1, \ldots, r_A\}$ from $|\Psi(\mathbf{R})|^2$.
      \item Propose moves: single-particle displacement; tune acceptance rate $\sim$50\%.
    \end{itemize}
  \item \textbf{Compute local energy.}
    \begin{equation*}
      E_L(\mathbf{R}) = \frac{H \Psi(\mathbf{R})}{\Psi(\mathbf{R})}
    \end{equation*}
  \item \textbf{Estimate the variational energy} $E_V = \langle E_L \rangle$ and its statistical error.
\end{enumerate}

\subsection{Phase 4: Stochastic Reconfiguration Optimization}
\begin{enumerate}[label=\textbf{4.\arabic*}]
  \item \textbf{Compute the overlap matrix $S$} and energy gradient $f$:
    \begin{align*}
      S_{ij} &= \langle O_i O_j \rangle - \langle O_i \rangle \langle O_j \rangle \\
      f_i &= \langle E_L O_i \rangle - \langle E_L \rangle \langle O_i \rangle
    \end{align*}
    where $O_i = \partial \ln \Psi / \partial \theta_i$.
  \item \textbf{Update parameters}: $\delta \theta = -\eta S^{-1} f$ (with regularization).
  \item \textbf{Iterate} until energy converges (typically $10^3$--$10^4$ iterations).
\end{enumerate}

\subsection{Phase 5: Results and Validation}
\begin{enumerate}[label=\textbf{5.\arabic*}]
  \item Compute ground-state energies for $^2$H, $^3$H, $^3$He, $^4$He.
  \item Compare to exact diagonalization (for small systems) and published VMC/GFMC results.
  \item Reproduce convergence curves (energy vs.\ iteration) from the paper.
  \item Study dependence on network architecture (depth, width).
\end{enumerate}

%---------------------------------------------------------------------
\section{Datasets}
%---------------------------------------------------------------------

\begin{longtable}{p{4.5cm} p{3.5cm} p{6cm}}
\toprule
\textbf{Dataset / Source} & \textbf{Type} & \textbf{Location / Notes} \\
\midrule
\endfirsthead
\toprule
\textbf{Dataset / Source} & \textbf{Type} & \textbf{Location / Notes} \\
\midrule
\endhead
NN potential parameters & Published tables & Argonne AV6$'$/AV8$'$: Wiringa et al.\ (1995, 2002); Minnesota: Thompson et al.\ (1977) \\
\midrule
Reference energies & Published benchmarks & GFMC results from Carlson et al., Pudliner et al. \\
\midrule
All other data & Synthetic (MC samples) & Generated by the VMC code \\
\bottomrule
\end{longtable}

%---------------------------------------------------------------------
\section{Models, Tools, and Software}
%---------------------------------------------------------------------

\begin{longtable}{p{4.5cm} p{2.5cm} p{6.5cm}}
\toprule
\textbf{Tool / Library} & \textbf{Version} & \textbf{Purpose / URL} \\
\midrule
\endfirsthead
\toprule
\textbf{Tool / Library} & \textbf{Version} & \textbf{Purpose / URL} \\
\midrule
\endhead
PyTorch & $\geq 1.10$ & Neural network, automatic differentiation \newline \url{https://pytorch.org/} \\
\midrule
JAX & latest & Alternative: JIT-compiled autodiff \newline \url{https://github.com/google/jax} \\
\midrule
NumPy & $\geq 1.21$ & Array operations \\
\midrule
SciPy & $\geq 1.7$ & Linear algebra (SR matrix inversion) \\
\midrule
Matplotlib & $\geq 3.4$ & Energy convergence and wave function plots \\
\midrule
NetKet & latest & Optional: VMC framework with SR built in \newline \url{https://www.netket.org/} \\
\bottomrule
\end{longtable}

\subsection{Hardware Requirements}
\begin{itemize}
  \item \textbf{$A \leq 4$}: Single GPU or multi-core CPU; runs in hours to days.
  \item \textbf{Memory}: Modest ($\sim$8\,GB RAM).
\end{itemize}

\end{document}
